\documentclass[11pt,a4paper]{article}

\usepackage{fontspec}
\usepackage{geometry,microtype,fancyhdr}
\usepackage{amsmath,amssymb}
\usepackage{booktabs,longtable,tabularx}
\usepackage{graphicx}
\usepackage{caption,subcaption}
\usepackage{enumitem}
\usepackage{listings}
\usepackage{titlesec}
\setlist{nosep}

\setmainfont[
    Path=../fonts/,
    UprightFont={Merriweather-VariableFont_opsz,wdth,wght.ttf},
    ItalicFont={Merriweather-Italic-VariableFont_opsz,wdth,wght.ttf}
]{Merriweather}
\setsansfont[
    Path=../fonts/,
    UprightFont={Inter-VariableFont_opsz,wght.ttf},
    ItalicFont={Inter-Italic-VariableFont_opsz,wght.ttf}
]{Inter}
\setmonofont[
    Path=../fonts/,
    UprightFont={GeistMono-VariableFont_wght.ttf},
    ItalicFont={GeistMono-Italic-VariableFont_wght.ttf}
]{Geist Mono}
\newfontfamily\displayfont[
    Path=../fonts/,
    UprightFont={PlayfairDisplay-VariableFont_wght.ttf},
    ItalicFont={PlayfairDisplay-Italic-VariableFont_wght.ttf}
]{Playfair Display}

\usepackage[shore]{../lrssntheme}
\usepackage{hyperref}
\usepackage{cleveref}
\hypersetup{
    colorlinks=true,
    linkcolor=LrssnPrimaryText,
    citecolor=LrssnPrimaryText,
    urlcolor=LrssnPrimaryText,
    pdftitle={Scientific Paper Title},
    pdfauthor={Author Name}
}

\geometry{left=22mm,right=22mm,top=24mm,bottom=24mm,headheight=22pt}
\setlength{\parindent}{0pt}
\setlength{\parskip}{0.65\baselineskip}

\titleformat{\section}{\displayfont\Large\bfseries\color{LrssnPrimaryText}}
{\thesection}{0.7em}{}
\titleformat{\subsection}{\sffamily\large\bfseries\color{LrssnSecondaryText}}
{\thesubsection}{0.7em}{}
\captionsetup{font=small,labelfont={bf,color=LrssnPrimaryText},
              textfont={color=LrssnMuted}}

\pagestyle{fancy}
\fancyhf{}
\fancyhead[L]{\small\color{LrssnMuted}Short title}
\fancyhead[R]{\small\color{LrssnMuted}\nouppercase{\leftmark}}
\fancyfoot[C]{\small\color{LrssnMuted}\thepage}

\makeatletter
\renewcommand{\maketitle}{%
  \begin{center}
    {\displayfont\LARGE\bfseries\color{LrssnPrimaryText}\@title\par}
    \vspace{0.7em}
    {\large\@author\par}
    {\small\color{LrssnMuted}\@date\par}
  \end{center}
  \vspace{1em}
}
\makeatother

\title{Scientific Paper Title}
\author{Author Name\\
        Department, Institution\\
        \texttt{author@example.org}}
\date{\today}

\begin{document}
\maketitle

\begin{abstract}
\noindent
State the research problem, approach, principal result, and significance.
Keep the abstract to a single focused paragraph.
\end{abstract}

\textbf{Keywords:} first keyword; second keyword; third keyword

\section{Introduction}
\label{sec:introduction}

Introduce the field, explain the problem, review its context, and state the
research question or hypothesis.

\section{Methods}
\label{sec:methods}

Describe data, materials, variables, experimental setup, and analysis methods
in sufficient detail for replication.

\subsection{Analysis}
Define equations and analytical procedures here.
\[
    \hat{y}_i = f(x_i;\theta)
\]

\section{Results}
\label{sec:results}

\begin{table}[htbp]
\centering
\caption{Example experimental results}
\begin{tabular}{lrrr}
\toprule
\textbf{Condition} & \textbf{Mean} & \textbf{SD} & \textbf{$n$} \\
\midrule
Control & 10.2 & 1.4 & 30 \\
Treatment & 12.8 & 1.7 & 30 \\
\bottomrule
\end{tabular}
\label{tab:results}
\end{table}

Report the findings first, then refer to \cref{tab:results} or
\cref{sec:methods} when relevant.

\section{Discussion}
\label{sec:discussion}

Interpret the results, relate them to previous work, discuss limitations,
and explain practical or theoretical significance.

\section{Conclusion}
Give a concise answer to the research question and identify future work.

\begin{thebibliography}{9}
\bibitem{example}
Surname, A. (2026). \emph{Example Scientific Work}.
Example Journal, 1(1), 1--10.
\end{thebibliography}

\end{document}
